\Needspace{6\baselineskip}
\Needspace{6\baselineskip}
\section{Terminal publication of the unique continuum through \(5\) non-permutable readers}
\Needspace{6\baselineskip}
\subsection{Factorization of the common enriched state into \(5\) terminal publications}
\Needspace{6\baselineskip}
\subsubsection{Causal dependence of terminal realizations with respect to the fully constructed genealogical continuum}
The five problems do not summon five distinct mathematical structures. Nor do they require the construction of five independent paths that, after advancing separately, must be artificially reunited in a common conclusion. That picture belongs to an earlier stage of the exposition. The current point of departure is stronger: before the names Riemann, Navier–Stokes, Yang–Mills, Hodge, or Birch–Swinnerton–Dyer appear, there already exists a single enriched mathematical object whose internal operations simultaneously produce the features required for the five closures.\par
For this reason, this account does not begin with problems. It begins with the construction of the continuum.\par
Triadic Modular Holography does not receive the continuum as a previously given background. Nor does it begin with the real line, a classical function space, an algebraic variety, a gauge connection, or an elliptic curve. It begins with a finite arithmetic endowed with orientation, memory, and the capacity for prolongation. APP provides the positional and arithmetic substrate; TRIT converts each operation into an oriented transition; TPK preserves the complete genealogy of those transitions and permits them to be prolonged without erasing the path that produced them.\par
The decisive step is that the output of an operation ceases to be merely a value. Each emission carries with it its arithmetic sheet, residue, quotient, orientation, carry, memory, boundary, route, multisection, survival, terminal object, and reader. The result is no longer an isolated point, but a state that knows where it comes from, how it can continue, and which relations it must preserve when the scale changes.\par
APP, TRIT, and TPK do not form three overlapping explanations. They form a single causal sequence:\par
\[
\mathrm{APP}
\longrightarrow
\mathrm{TRIT}
\longrightarrow
\mathrm{TPK}
\longrightarrow
\text{unique enriched state}.
\]
APP fixes the common arithmetic support. TRIT introduces the effective orientation of each transition, distinguishing advance, threshold, and future opening. TPK brings together the visible information and the transported information: not only what an operation yields, but also the carry it shifts, the boundary it modifies, the memory it preserves, and the extensions it permits.\par
The nonadic connection does not return the system to an empty state. After one complete turn, the phase reappears, but the history is not erased. The return \(9\to1\) is phase return, not state reset. Every turn adds depth, memory, and compatibility. Continuity is therefore obtained neither by repeating nine filters nor by traversing nine independent channels. It is obtained by prolonging a single complete connection whose nine segments are internal charts of the same holonomy.\par
\Needspace{6\baselineskip}
\subsubsection{Correlative emergence of \texorpdfstring{\(5\)}{5} non-exhaustive intrinsic actions}
A common operation acts on every enriched state. Five domains, five selections, or five copies of the state do not exist first. The same emission is read simultaneously under five internal characteristics denoted by the discriminants \(\mathrm{sp}\), \(\mathrm{sol}\), \(\mathrm{gau}\), \(\mathrm{coh}\), and \(\mathrm{det}\).\par
Those symbols do not name branches or exhaust the properties of the continuum. They name
\(5\) intrinsic actions of the same object, whose correlation is decisive
for the terminal realizations of this Part.\par
The spectral–polar characteristic, \(sp\), records the oriented publication of the state, its complement, its defect, its terminal memory, and the relation between the two mirror readings of the same form.\par
The solenoidal characteristic, \(sol\), records how a transition preserves state and, at the same time, distributes dissipation, advection, carry, memory, microstructure, and shell without overcounting them.\par
The gauge characteristic, \(gau\), records local equivalence, incidence, holonomy, curvature, and compatible passage to the physical quotient.\par
The cohomological characteristic, \(coh\), records the extensions between levels, the relative obstructions, the compatibility of the lifts, and the survival of a class through every refinement.\par
The determinantal characteristic, \(det\), records the finite-global incidence, the filtration pages, the local complexes, their determinant lines, and the transport of an arithmetic basis up to an analytic publication.\par
But none of these properties exists in isolation. All act on the same emissions and preserve the same genealogical thickness. The fiber carrying the spectral characteristic is the same enriched fiber whose orientation participates in the solenoidal balance, whose boundary supports the gauge reading, whose prolongation sustains the cohomological extension, and whose incidence enters the determinant line.\par
Five objects are not generated and then packaged. A single object is generated that inseparably possesses those five capacities.\par
A posterior view is therefore not a projection that creates a branch. It is solely a way of displaying, for a specific application, a characteristic that was already operating within the whole. The view does not retrospectively select favorable states or preserve the input as a first coordinate in order to recover it tautologically later. It makes visible a structure produced and preserved throughout the entire genealogy.\par
\Needspace{6\baselineskip}
\subsubsection{Causal precedence of the already constructed continuum as the limit of the compatible prolongation}
At each depth there exists a finite state that is nevertheless complete at its level. Extensions between levels cannot arbitrarily alter prior information. They must preserve the fibers, extension relations, numbers, orientation, carry, memory, boundary, route, multisection, survival, terminal object, and readers.\par
Naturality requires that operating and then truncating produce the same result as truncating and then operating. Owing to that compatibility, finite histories do not remain disconnected approximations. They form an inverse system. The continuum appears when the complete family of compatible histories is taken, not when an infinite collection of points is postulated in advance.\par
The HMT continuum is therefore coherent memory of paths. A point without genealogy would be an ablation: it would preserve the visible value while eliminating precisely what makes it possible to prove how that value survives every refinement.\par
Only after that compatibility has been closed does the complete discrete structure of the continuum appear. Only then can an Archimedean realization be published. The real line is not the raw material of the construction; it is an output of one of its terminal readers. The same holds for the classical domains of the five problems: they appear when the continuum has already been constructed and a subsequent interface translates one of its intrinsic features into the language of the corresponding target.\par
The correct causal sequence is, therefore:\par
\[
\begin{aligned}
\mathrm{APP}
&\longrightarrow \mathrm{TRIT}
\longrightarrow \mathrm{TPK}
\longrightarrow \text{unique enriched state}\\
&\longrightarrow 04\mathrm b1
\longrightarrow 04\mathrm b2
\longrightarrow 04\mathrm b3\\
&\longrightarrow \text{correlative intrinsic actions}
\longrightarrow \text{terminal and limit}
\longrightarrow \mathfrak C_{\mathrm{cont}}^{\mathrm{disc}}.
\end{aligned}
\]
Then, and not before, enter the specific entanglers:\par
\[
\mathfrak C_{\mathrm{cont}}^{\mathrm{disc}}
\longrightarrow
\begin{cases}
\text{explicit Weil form},\\
\text{three-dimensional incompressible evolution},\\
\text{quantum gauge theory},\\
\text{publication of Hodge classes},\\
\text{BSD determinantal comparison}.
\end{cases}
\]
These five terminal translations do not define the structure they use. The classical problem supplies its data afterward: a test function in Riemann, initial data and a force in Navier--Stokes, a simple compact group in Yang--Mills, a variety and a class in Hodge, or an elliptic curve in BSD. None of those data participates in the prior generation of the continuum.\par
The causal direction is thereby protected. The target does not select the state that is to solve it. The state exists before the target; the intertwiner merely proves that the internal structure and the classical formulation express the same operation in two languages.\par
\Needspace{6\baselineskip}
\subsection{1st resolution: Riemann and positivity as a constructed complement}
The Riemann hypothesis is usually presented as a statement concerning the location of the nontrivial zeros of the zeta function. From that formulation, the problem appears to call for inspecting a given population of zeros and proving that all of them align on the critical line.\par
The HMT reading reverses that perspective. It does not begin by asking
where the zeros lie. It asks what structure simultaneously generates both
sides of the explicit formula and what prevents their difference from
acquiring a negative sign.\par
The spectral–polar characteristic of the continuum provides a double publication of the same state. One orientation collects the positive channel; the other collects its mirror image. Both contain, under a common regulator, the prime contribution, the gamma–scalar contribution, and the polar contribution. They are not three distinct proofs or three sums appended at the end. They are three rows of the same colligation constructed from the starting data.\par
The essential point is that this colligation is constructed before invoking the positivity that it must demonstrate. Its input, outputs, loss, and terminal state are typed within the same enriched space. The identity linking the two publications is not introduced as a desired equality between classical formulas; it arises from the internal dynamics of the connection.\par
Between the input publication and the mirror publication, a residual output
may appear at finite depth. The nonadic structure compels it to satisfy an
antimirror recurrence. At each complete turn, the next residual is
transported isometrically but multiplied by the factor\par
\[
q=\frac{5}{9}.
\]
If \(\Omega_n\) represents that defect at depth \(n\), the recurrence has the form\par
\[
\Omega_n=q\,V_n\Omega_{n+1},
\]
with \(V_n\) isometric and the family coinductively bounded. Upon iterating \(N\) times,\par
\[
\|\Omega_n\|
\le q^N\sup_{m\ge n}\|\Omega_m\|.
\]
As \(0<q<1\), the right-hand term tends to zero. Therefore, the residual does not vanish because it was omitted or because an identification between the two publications was assumed. It vanishes because the complete recurrence contracts it at every compatible depth. The specular output then coincides with the constructed negative publication.\par
Once that link has been discharged, the isometry constructed from the starting data yields the energy identity. At finite depth, the norm of the positive publication decomposes into three parts: the negative publication, the accumulated loss, and the terminal state. The terminal is explicitly preserved. It must be preserved through passage to the limit because it contains the memory of truncation and certifies that telescoping prolongs the finite identity through a proved convergence.\par
The finite norm decomposition satisfies the exact identity:\par
\[
\|\text{positive publication}\|^2
=
\|\text{negative publication}\|^2
+
\|\text{losses up to }N\|^2
+
q^{2N}\|\text{terminal}\|^2.
\]
Each summand has a concrete origin. The negative publication is the already identified spectral output. The losses gather exactly the complement that was not published by that orientation. The terminal remembers what still remains beyond the cut.\par
Only after the complete identity has been written is \(N\to\infty\) imposed. Then, and only then, does the terminal term disappear because \(q^{2N}\to0\). The difference between the two publications is converted into the square of a loss operator:\par
\[
\mathcal I_{\mathrm{GW}}
=
E_R^{*}E_R.
\]
This is the Guinand--Weil form. Its positivity is no longer a conjecture added to the construction:\par
\[
\langle f,\mathcal I_{\mathrm{GW}}f\rangle
=
\|E_Rf\|^2
\ge0.
\]
The HMT structure does not obtain this inequality by concealing the negative part, restricting the domain to functions that already satisfy it, or introducing a projector whose existence is equivalent to the conclusion. It constructs the complement and proves that the form is exactly its square.\par
The common publication is then materialized. A single encoder preserves the prime, gamma–scalar, and polar channels. An internal projector produces one of the two moments; its complement produces the other. The difference between the moments is not an external comparison between similar formulas: it is the norm of the orthogonal component constructed by the colligation.\par
As the window is enlarged, individual columns are not forced to satisfy false isometries between distinct truncations. What is preserved cofinally is the form, the reachable complement, and naturality of its publication. The increments entering when the cut is enlarged appear in both orientations with the same neutral contribution and do not alter the governing difference. This yields a compatible family of positive forms on a separating class.\par
The final transition is the reversible Weil criterion. Classical theory states that positivity of the explicit form, on the appropriate class of test functions, is equivalent to the location of all nontrivial zeros on the critical line. HMT supplies precisely that positivity, but supplies it as a consequence of an operator identity constructed within the continuum.\par
\Needspace{5\baselineskip}
Por tanto,\par
\[
\operatorname{Re}\rho=\frac12
\]
for every nontrivial zero \(\rho\) of the zeta function.\par
The critical line appears here in two concordant ways. Geometrically, it is the fixed locus of the involution that exchanges the two spectral orientations. Operatorially, it is the only position compatible with the positivity of the publication complement. Symmetry explains why the central line is distinguished; the quadratic identity proves why the zeros cannot leave that line.\par
The resolution does not consist in observing that the functional equation has a symmetry. It consists in constructing the object whose defect measures the breaking of that symmetry, proving that the defect vanishes coinductively, and converting the remaining difference into a norm. It is the combination of genealogy, terminal object, contraction, common publication, and positivity that closes the problem.\par
Thus the first of the five terminal resolutions is established. Not as an independent branch, but as the classical translation of the spectral–polar characteristic that the unique continuum possessed from its construction.\par

\Needspace{6\baselineskip}
\section{Dynamic closures of the continuum: solenoidal transport and gauge quotient}
\Needspace{6\baselineskip}
\subsection{Dynamics and curvature of the genealogical continuum}
The first resolution showed how the spectral--polar characteristic of the continuum is published as Weil positivity. The next two resolutions interrogate the same state through two other consubstantial aspects: its capacity to transport nonlinear dynamics without losing energy, memory, or regularity, and its capacity to support a local gauge theory whose vacuum is separated by a positive spectral gap.\par
Navier–Stokes and Yang–Mills do not receive two new structures. They receive two distinct intertwiners from the same already constructed continuum. In the first case, the intertwiner publishes the state as an incompressible velocity field with viscosity, advection, and pressure. In the second, it publishes the state as a refinable gauge network with holonomy, curvature, Euclidean measure, and a physical Hamiltonian.\par
The difference between the two problems does not break the unity of the construction. In Navier--Stokes, the critical point is to prove that no nonlinear transfer disappears outside the balance and that the regularity obtained is uniform as the cutoffs are removed. In Yang--Mills, the critical point is to prove that the finite gap belongs to the same gauge theory that reaches the continuum, survives the physical quotient, and acts in the curvature sector rather than in an auxiliary factor.\par
\Needspace{6\baselineskip}
\subsection{Global existence, regularity and uniqueness for Navier–Stokes through a uniform bound and compactness}
The three-dimensional Navier–Stokes problem does not arise from the mere presence of a nonlinear term. The difficulty arises because the classical publication of the velocity field tends to obscure where transfer between scales is stored. Viscosity is visible. Kinetic energy is visible. Advection also appears in the equation. Yet the complete accounting of the interaction among depth, memory, carry, microstructure, and shell is compressed into a single visible function.\par
That compression is precisely what the solenoidal character of the continuum avoids.\par
The enriched state does not retain only the field \(u\). It simultaneously retains the Galerkin field, the time orientation, the positive and signed memories, the carry between cuts, the microscopic component, the shell of new frequencies, the route, the evolution archive, and the terminal. The classical equation is a publication of that state; it is not its complete definition.\par
Given initial data \(u_0\), force \(f\), viscosity \(\nu>0\), and regularity \(s>5/2\), the Galerkin projection produces\par
\[
\partial_tu_M+\nu A_Mu_M+B_M(u_M,u_M)=f_M.
\]
But \(M\) does not select a favourable solution. It merely fixes a finite resolution in which all operations can be written exactly. The objective is to construct an estimate that does not depend on \(M\) or on the nonadic depth \(J\). Only a bound with that uniformity can pass to the continuum.\par
\Needspace{6\baselineskip}
\subsubsection{Decomposition of visible memory and genealogical memory in the evolutionary equation}
The nonlinear interaction has a sign. At certain times it delivers energy to one scale; at others it withdraws it. If only its signed value is preserved, the total magnitude transferred may be lost. If only its absolute value is preserved, orientation is lost.\par
The structure preserves both.\par
One memory \(D_M\) publishes the sign of the advective transfer. Another memory \(H_M\) preserves its positive magnitude. Their separation makes it possible to record simultaneously\par
\[
\dot D_M=\mathscr B_{s,M},
\qquad
\dot H_M=\frac{1+r}{r-1}|\mathscr B_{s,M}|,
\]
where \(r=\pi_{\rm HMT}/\varphi_{\rm HMT}^2>1\).\par
This adds no fictitious energy. It unfolds two distinct aspects of a single transfer: orientation and cost. Merging them would destroy information; adding them without orthogonality would duplicate energy. The construction will make explicit that neither of those two deformations occurs.\par
\Needspace{6\baselineskip}
\subsubsection{Independence of the initial data with respect to the subsequent genealogical history}
The most serious objection to any operatorial closure of Navier–Stokes would be that the uniform bound had been introduced retrospectively: first, a regular solution would be chosen; its history would then be encoded; and finally it would be announced that the code controls that same solution.\par
The construction prevents that circularity by fixing the source space before raising the evolution:\par
\[
\mathcal D_T
=
H^s_\sigma(\mathbb T^3)
\times
L^2(0,T;H^{s-1}_\sigma(\mathbb T^3)),
\qquad
d=(u_0,f).
\]
The root \(J_T^{\rm sol}d\) depends exclusively on \(u_0\) and \(f\). It does not receive as input \(u_M\), the future loss, the terminal object, or the norm that must be bounded.\par
Force is likewise not introduced as an artificial additive port. The Carleman equation shows that it acts in a mixed manner: every occurrence of \(f\) is combined with a preceding degree of the state. The lift therefore uses a mixed state–force creation. The chronological signature of \(f\) records all its ordered times, whereas the PC–Fock tower records all nonlinear degrees of the state.\par
The composition materially reproduces the Dyson series of the nonautonomous generator. Applying the normalized degree-one reader recovers exactly \(u_M(t)\). What is recovered is neither a symbolic approximation nor a variable resembling the solution: it is the solution of the Galerkin ODE, because both satisfy the same finite equation and the same initial datum.\par
The causal construction from the initial data therefore proceeds in the order\par
\[
(u_0,f)
\longrightarrow
J_T^{\rm sol}(u_0,f)
\longrightarrow
\text{state–force tower}
\longrightarrow
\text{Galerkin history}.
\]
The genealogical history is constructed as a later output of the data of departure and is excluded from its initial definition.\par
\Needspace{6\baselineskip}
\subsubsection{The scattering identity that preserves cross power}
Force and dissipation are organized by means of two waves:\par
\[
a_T(f)
=
\frac{1}{2\sqrt\nu}A^{-1/2}\Lambda^sP_\sigma f,
\qquad
z_M(u)=\sqrt\nu A_M^{1/2}\Lambda^su,
\]
and the residual wave\par
\[
b_M=z_M-P_Ma_T(f).
\]
The applied power is not replaced by the norm of the force. It is preserved exactly as a cross term:\par
\[
\mathscr F_{s,M}
=
2\operatorname{Re}\langle z_M,P_Ma_T(f)\rangle.
\]
Upon completing the square, the exact balance is obtained.\par
\[
\dot G_{M,j,T}
+\|b_M\|^2
+\lvert\mathscr B_{s,M}\rvert
=
\|P_Ma_T(f)\|^2.
\]
This is one of the points that make the closure work. If \(b=z-a\) were replaced by \(z\), the cross term representing the work of the force would disappear and the identity would be false. The structure does not control the nonlinearity by discarding it: it turns it into an oriented positive channel and also preserves the exact interference between force and dissipation.\par
\Needspace{6\baselineskip}
\subsubsection{Factorization of \(6\) losses through a unique operator without spurious multiplicity}
The complete history distinguishes six outputs:\par
\begin{itemize}
\item dissipation;
\item oriented advection;
\item carry between cuts;
\item microscopic memory;
\item shell of new frequencies;
\item non-inherited memory.
\end{itemize}
They are not six energies added to a previous balance. They are six orthogonal rows of a single terminal loss operator. Parseval's identity establishes\par
\[
\|L_{M,j,T}^{\rm term}x\|^2
=
\sum_{\mathfrak a}
\|L_{M,j,T}^{\mathfrak a}x\|^2.
\]
The orthogonality is not merely nominal. The depths occupy disjoint temporal intervals; carry and shell inhabit distinct spectral ranges; the microstructure belongs to the complement \(Q_{\rm micro}\); the new memory is separated from the inherited memory; and the two sheets of the advection cannot be positive simultaneously at the same point.\par
For this reason, the accounting does not duplicate the interaction. The sum of the six norms is exactly the norm of the common operator. If one of them were not a projection of that operator, Parseval would fail before the bound was obtained. The formalism itself thus contains its falsifier.\par
\Needspace{6\baselineskip}
\subsubsection{Persistence of the terminal object under passage to the limit}
Every truncation preserves something that has not yet been published. That remainder is the terminal object. It contains the visible energy at the end of the interval and the oriented memory that remains between the observed time and the horizon \(T\).\par
The terminal column is constructed first. Its Gram matrix is then formed. A desired metric is not defined and an operator subsequently invented to realize it. The column materially contains the terminal output and all transported losses:\par
\[
\mathbf G_{M,j;J,T}^{\rm term}
=
\begin{bmatrix}
\text{transported terminal}\\
\text{loss}_{J-1}\\
\vdots\\
\text{loss}_{j}
\end{bmatrix}.
\]
His block recursion produces the Stein identity.\par
\[
U_{M,j,T}
=
\Gamma_{M,j,T}^{*}
U_{M,j+1,T}
\Gamma_{M,j,T}
+
(L_{M,j,T}^{\rm term})^{*}
L_{M,j,T}^{\rm term}.
\]
This identity does not retroactively define the norm of the state. It is obtained by multiplying an already constructed column by its adjoint. For the same reason, global telescopy is an identity among existing operators, not a stability hypothesis.\par
Future memory is not introduced into the initial root either. It resides in the terminal output. Upon taking the integral estimate over \(\tau=T\), its temporal support is emptied exactly. It is not erased for convenience.\par
\Needspace{6\baselineskip}
\subsubsection{Uniform a priori bound and Aubin–Lions compactness}
The root, the level encoder, the column, and the common isometry satisfy the factorization derived from the data of departure\par
\[
\mathcal G_{M,J,T}^{\rm sol}
\Lambda_{M,T}
J_T^{\rm sol}d
=
I_{M,J}J_T^{\rm sol}d,
\qquad
I_{M,J}^{*}I_{M,J}=I.
\]
The left-hand side is the complete history at the cutoff \((M,J)\). The right-hand side depends on the unique root constructed from the data. Since \(I_{M,J}\) is isometric, the norm of the history is dominated by the source norm with a constant independent of \(M\) and \(J\).\par
This is the step that resolves the difficulty. Uniformity does not follow from a finite-dimensional estimate whose constant grows with the cutoff. Nor does it follow from defining the history space through the very loss one seeks to control. The root exists before the cutoff, and all cutoffs are isometric realizations of it.\par
The terminal column translates that structural bound into the physical quantities:\par
\[
\begin{aligned}
&\sup_M\sup_{0\le t\le T}\|u_M(t)\|_{H^s}^2
+\nu\sup_M\int_0^T\|u_M(t)\|_{H^{s+1}}^2\,dt\\
&\quad
+\sup_M\int_0^T|\mathscr B_{s,M}(u_M(t))|\,dt
+\text{carry}+\text{micro}+\text{shell}+\text{memory}
\le C_T(u_0,f).
\end{aligned}
\]
The constant is finite for every \(T<\infty\) and does not depend on \(M\) or \(J\). The inequality controls not only the visible field. It simultaneously controls the transported work and the sectors that a classical publication usually forgets.\par
\Needspace{6\baselineskip}
\subsubsection{Sign Signature Obstruction in Kalman–Yakubovich–Popov Factorization}
Making the global closure depend on a KYP factorization with an incorrect sign introduces a positive square and purports to recover it afterward with a negative sign. That equality is algebraically false.\par
The proof excludes that identity and directly constructs the necessary uniform estimate.\par
The retained KYP chart is a correct finite identity under its own conditions. It serves as an auxiliary control for each cutoff, but does not provide uniformity when \(M\) and \(J\) are removed. The global proof does not use a bound on \(Q^{-1}\), does not take the limit within KYP, and does not depend on that sign.\par
Therefore, the obstruction from the incorrectly signed KYP block does not affect the proof. Uniformity is established by the direct construction from the data with its terminal column and its Stein identity.\par
\Needspace{6\baselineskip}
\subsubsection{Convergence from the uniform bound to the global classical solution}
The previous bound provides\par
\[
u_M
\quad\text{bounded in}\quad
L^\infty(0,T;H^s)
\cap L^2(0,T;H^{s+1}).
\]
Since \(s>5/2\), the Sobolev product controls the nonlinearity:\par
\[
B:H^s_\sigma\times H^s_\sigma
\longrightarrow H^{s-1}_\sigma.
\]
The Galerkin equation further provides a uniform bound for \(\partial_tu_M\) in \(L^2(0,T;H^{s-1})\). The inclusions\par
\[
H^{s+1}\Subset H^s\hookrightarrow H^{s-1}
\]
allow the application of Aubin--Lions. Strong convergence in \(L^2H^s\) is obtained, sufficient to transport the bilinear term:\par
\[
B_M(u_M,u_M)\longrightarrow B(u,u)
\quad\text{in}\quad
L^1(0,T;H^{s-1}).
\]
The limit satisfies the incompressible equation. Pressure is not introduced as another hidden datum; it is reconstructed afterwards through\par
\[
-\Delta p=\partial_i\partial_j(u_i u_j),
\qquad
\int_{\mathbb T^3}p=0.
\]
Uniqueness is obtained from the difference of two solutions. An estimate in
\(H^{s-1}\), followed by Grönwall's inequality, forces the difference to
vanish. This eliminates dependence on the subsequence extracted by
compactness.\par
The construction holds for every finite horizon. The solutions obtained in \([0,n]\) and \([0,n+1]\) agree on their overlap by uniqueness. They therefore glue into a single solution defined over \([0,\infty)\).\par
Smoothness is obtained throughout the scale \(s+n\) using the same structural readers; uniqueness identifies all limits and produces a global smooth solution for every smooth periodic datum. The required bound is constructed on each finite horizon and at every Sobolev index from the complete solenoidal datum, as proved in the source development \cref{thm:pdf2-ns-global}.\par
Finally, the zero-mean condition is a normalization of the solenoidal sector. The time-dependent Galilean transformation separates the mean and bijectively transports every smooth periodic datum to the zero-mean sector, preserving divergence, smoothness and Sobolev norms.\par
\Needspace{6\baselineskip}
\subsubsection{Global existence, regularity and uniqueness for Navier--Stokes}
The solenoidal characteristic of the continuum, expressed through its specific intertwiner, produces a unique global smooth solution of the three-dimensional periodic Navier--Stokes equations for every smooth periodic datum. The Leray--Hopf construction is contained as an energy stage within that strong conclusion.\par
The defense of this conclusion does not rest on saying that a structure
capable of solving the problem exists. It rests on five material facts:\par
the root uses only \((u_0,f)\); the nonlinearity is reproduced by mixed creation; the six losses are orthogonal and are not added twice; the terminal remains until telescopy; and the uniform bound precedes the classical compactness argument.\par
Aubin–Lions, pressure, Grönwall, and gluing are not the source of regularity. They are the final classical recognition of a regularity that the structure has already prevented from being lost.\par
\Needspace{6\baselineskip}
\subsection{3rd solution: Yang–Mills and the curvature sector gap}
The Yang–Mills problem requires two inseparable things: constructing a nontrivial quantum theory over \(\mathbb R^4\) for every compact, connected, and simple group of the considered domain, and proving that its Hamiltonian has a positive mass gap.\par
It is not enough to exhibit a finite operator with positive spectrum. Nor is it enough to write down a Wilson action, a product measure, or an observable with a nonzero eigenvalue. The gap must belong to the same gauge theory that possesses locality, reflection positivity, Euclidean covariance, Wightman reconstruction, and physical curvature.\par
The gauge characteristic of the continuum already preserves the refinable network, the measure, the four reflections, the holonomy, color, incidence, the fusion carry, the route, memory, the terminal object, and the curvature reader. The group \(G\) and the coupling \(g\) enter afterward through the specific intertwiner. They do not retrospectively select the genealogy.\par
\Needspace{6\baselineskip}
\subsubsection{Non-equivalence between the simple fiber and the product dynamics}
The first defensive distinction separates two generators that an old notation could confuse.\par
The semigroup of a alone fiber has only the levels \(0\) and \(3\kappa_{\rm av}\). The generator product complete integra, in change, all the coordinates independent of the state:\par
\[
D_m^{\rm prod}
=
3\kappa_{\rm av}
\sum_{\iota\in\mathcal I_m}(I-E_{m,\iota}),
\]
where the expectations \(E_{m,\iota}\) are orthogonal and commuting.\par
Its spectrum contains\par
\[
0,\quad
3\kappa_{\rm av},\quad
6\kappa_{\rm av},\quad
9\kappa_{\rm av},\ldots
\]
according to the number of nonconstant coordinates. Replacing this generator with that of a simple fiber would erase the Hoeffding multiplicities and would be immediately falsifiable: a vector depending on two factors must have energy \(6\kappa_{\rm av}\), not \(3\kappa_{\rm av}\).\par
The gap will be proved for the product generator and will then be transported to the physical sector. It is not obtained by confusing two different dynamics.\par
\Needspace{6\baselineskip}
\subsubsection{Absence of hidden coordinates and exclusion of spurious null modes}
The complete state is factorized via an exhaustive index. It appears only once:\par
\begin{itemize}
\item frames;
\item links;
\item incidents;
\item marks;
\item refinement bridges;
\item new innovations at each level.
\end{itemize}
The infinite subtails are recorded by means of their scalar cylinders. No genealogical coordinate remains outside the family of expectations.\par
The Hoeffding decomposition writes the space as an orthogonal sum of sectors \(\mathcal H_{m,S}\). On each one,\par
\[
D_m^{\rm prod}u_S
=
3\kappa_{\rm av}|S|u_S.
\]
The sector \(S=\varnothing\) is the intersection of the ranges of all expectations and consists exactly of the constants:\par
\[
\bigcap_{\iota\in\mathcal I_m}
\operatorname{Ran}E_{m,\iota}
=
\mathbb C\Omega_m.
\]
Therefore, the vacuum is simple, and every vector orthogonal to the vacuum has energy at least \(3\kappa_{\rm av}\).\par
This excludes the objection that the gap arises from failure to enumerate a coordinate. If a factor were missing, there would be nonconstant functions invisible to all expectations, and additional zero modes would appear. The exhaustiveness of the index prevents exactly that failure.\par
\Needspace{6\baselineskip}
\subsubsection{Observable realization of the spectral gap}
A lower bound does not suffice to fix the exact value of the gap. A witness that attains the threshold is required.\par
In each incidence collar, the variable\par
\[
W_c(q)
=
\mathbf1_{\{q_1+\cdots+q_6=0\}}
-\frac13
\]
has zero mean and satisfies\par
\[
D_m^{\rm prod}W_c
=
3\kappa_{\rm av}W_c.
\]
It must still be proved, however, that \(W_c\) belongs to the physical gauge sector.\par
The incidence action is written through the matrix\par
\[
(B^+x)_{ij}=x_i+x_j.
\]
Each component of \(x\in\mathbb F_3^4\) appears three times when the six incidences are summed. In \(\mathbb F_3\),\par
\[
\sum_{i<j}(B^+x)_{ij}
=
3\sum_i x_i
=
0.
\]
Therefore, the condition defining \(W_c\) is gauge invariant. Averaging over the full gauge group preserves \(W_c\); it does not remove it from the physical space. Thus, \(W_c\) is a physical eigenvector, and the value \(3\kappa_{\rm av}\) is not merely a bound: it is the first spectral level attained.\par
\Needspace{6\baselineskip}
\subsubsection{Construction of the measure prior to the estimation of the gap}
A coarse edge is refined into nine ordered increments. The fine measure is not contrived by copying the coarse measure nine times. It is constructed by means of the heat bridge conditioned by the product of the nine increments. The times of the subsegments may differ and are subject only to the requirement that their sum equal the coarse time.\par
Convolution with the heat kernel ensures that the bridge is a probability. A triangular chart then separates the coarse variable from the new innovations. Projection onto the preceding level multiplies the nine increments and forgets only the innovations.\par
In this way,\par
\[
(\widehat\pi_{m+1,m})_*
\widehat\mu_{m+1}
=
\widehat\mu_m.
\]
Edges shared by several faces remain the same variable. They are not replaced with independent copies. Carry and memory prevent recounting at the fine level the face that was already present at the coarse level.\par
Over that projective family one first constructs the exact reader of \(F^2\), its positive root, the local action, and the coupling-dependent Gibbs density:\par
\[
d\widehat\mu_{m,g}^{\rm YM}
=
Z_{m,g}^{-1}
e^{-S_{m,g}}\,
d\widehat\mu_m^0.
\]
The Schwinger–Dyson and Ward equations are calculated for that measure. The gap is not used to choose \(g\), to define \(S_{m,g}\), or to reweight the law. It appears afterward, when studying the reversible generator associated with the same measure.\par
This is the defense against action--law--gap circularity: first the action and its measure are constructed; afterwards its spectrum is proved.\par
\Needspace{6\baselineskip}
\subsubsection{Descent of the generator to the quotient of caliber}
The complete space includes links, frames, incidences, color, orientation, route, memory, and terminal. The marginal that preserves only the links is a subsequent ablation. It cannot be used to deny an observable that resides in the incidence fiber.\par
The physical projection is obtained by averaging the action of the full gauge group. Since the generator integrates coordinates equivariantly,\par
\[
[\mathbb P_m^{\rm phys},D_m^{\rm prod}]=0.
\]
The semigroup therefore descends to the physical quotient. Its kernel is first defined on the complete state and only afterward on the gauge classes. One does not take an operator restricted to a fiber and declare it to be the kernel over the whole space.\par
The simple vacuum and the eigenvector \(W_c\) survive compression. The finite gap already belongs to the gauge-invariant space.\par
\Needspace{6\baselineskip}
\subsubsection{Uniqueness of the process reconstructed from \(4\) positivities by reflection}
The Euclidean theory must satisfy reflection positivity in the four directions. Four independent measures are not constructed. A single face measure, obtained from the same reversible kernel, factorizes each Osterwalder–Schrader form:\par
\[
\int
\overline{F(\Theta_{\nu}y)}F(y)\,d\Omega
=
\|\mathcal B_\nu F\|^2
\ge0,
\qquad \nu=1,\ldots,4.
\]
Positivity is not an added hypothesis. It is a squared norm produced by Markov and reversibility.\par
The marginals of opposite faces recover the same semigroup. Therefore, the four OS reconstructions have a single dynamical origin. The connectors between half-lattices will subsequently prove that they reconstruct the same space, the same vacuum, and the same Hamiltonian.\par
\Needspace{6\baselineskip}
\subsubsection{Persistence of the Spectral Gap under the Continuous Limit}
The refinement maps intertwine the Dirichlet forms, the semigroups, the vacuum, and the witness \(W_c\). Passage to the continuum is not based solely on observing that the gap has the same value at every level. In isolation, that equality would not prevent new states of arbitrarily small energy from appearing in the limit.\par
The construction controls cross-products through block intertwiners. The coarse coordinates are distributed over their descendants with exact normalization. The Gram identities prove that the smoothed publications form Cauchy families. The form, the vacuum, and the eigenvector thereby attain the limit within the same inductive system.\par
Se obtiene\par
\[
D_{\infty,g}^{\rm YM}
\succeq
3\kappa_{\rm av}(I-P_{\Omega_\infty}),
\]
and the limit witness satisfies\par
\[
D_{\infty,g}^{\rm YM}W_c
=
3\kappa_{\rm av}W_c.
\]
\Needspace{5\baselineskip}
Por tanto,\par
\[
\Delta_{\rm YM}^{\rm HMT}
=
3\kappa_{\rm av}>0.
\]
The inequality and the witness converge together. The gap is neither deduced from mere lower semicontinuity nor extrapolated from an isolated lattice.\par
\Needspace{6\baselineskip}
\subsubsection{Spectral witness membership in the curvature sector}
An objection would still remain if \(W_c\) were only a combinatorial excitation of an internal factor, decoupled from Yang–Mills curvature.\par
Incidence holonomy eliminates that possibility. The sum of the six incidences determines a local holonomy of order three; \(W_c\) is a class function of that holonomy. It therefore belongs to the local gauge-invariant algebra generated by the loops.\par
The reader \(F^2\) is polarized to produce a space of local curvatures. The refinement intertwiners transport those curvatures while preserving the \(\mathfrak g\) product. The cellular sums converge to a gauge-covariant distribution\par
\[
\mathbf F_{\mu\nu}
\in
\mathcal S'(\mathbb R^4;\mathfrak g)
\widehat\otimes
L^2(\widehat X_\infty,\widehat\mu_{\infty,g}^{\rm YM}).
\]
The cyclic sector generated by \(\mathbf F\) is reducing for the Hamiltonian and contains the witness \(W_c\). Consequently, the exact gap belongs to the physical curvature sector. It is not the spectrum of an auxiliary noise.\par
\Needspace{6\baselineskip}
\subsubsection{Euclidean locality \(E(4)\) and continuous Osterwalder–Schrader reconstruction}
The generator decomposes into local block circuits with uniformly bounded support. A nonadic coloring separates the necklaces that may be updated simultaneously. The product of the circuits preserves the Bianchi identity at every step, not only upon completion of a sweep.\par
Euclidean translations and rotations act by relabelling the lattice, frames, incidences, and holonomies. Strong continuity is not proved only for an abstract first chaos. It is first controlled on local observables, then on holonomies, plaquettes, and clover terms; orbit smoothing extends the action to a dense core.\par
The cofinal system produces a local net \(\mathfrak A_{\rm YM}(O)\), a vacuum \(\Omega_{\infty,g}\), a strong representation of \(E(4)\), and Schwinger functionals. The Markov stacks and OS transfers are identified with the same continuous semigroup. Reflection positivity, regularity, covariance, and exponential clustering permit Osterwalder--Schrader reconstruction.\par
The continuation produces local Wightman fields and a Poincaré representation with spectrum in the future cone. The four Euclidean reconstructions coinduce the same curvature Hamiltonian.\par
The gap also controls the cluster:\par
\[
|S^{\rm YM}(A\,\tau_rB)|
\le
e^{-3\kappa_{\rm av}r}
\|A\Omega_\infty\|
\|B\Omega_\infty\|.
\]
Thus, the positive spectral value has a continuous and local physical consequence; it does not remain confined to the finite calculation.\par
\Needspace{6\baselineskip}
\subsubsection{Reconstruction of the classical action as the output of the Euclidean system}
The governing finite density is the exact reader of \(F^2\). The clover does not replace it. It functions as a smooth chart for recognizing its limit.\par
For a smooth connection, the Magnus expansion of the holonomy of a plaquette starts with\par
\[
\log\operatorname{Hol}_{p_{\mu\nu}}
=
a_m^2F_{\mu\nu}
+\text{odd term}
+O(a_m^4).
\]
The average over the four clover orientations cancels the odd term. The Riemann sum gives\par
\[
S_{m,g}(A)
\longrightarrow
\frac1{4g^2}
\int_{\mathbb R^4}\|F_A(x)\|^2\,d^4x.
\]
Therefore, the constructed law is recognized as Yang–Mills because its finite action, holonomies, and curvature converge to the correct classical functional. A “Yang–Mills theory” is not defined as any theory possessing the desired gap.\par
The same measure satisfies Schwinger–Dyson and Ward. The relation between small loops and curvature, the Bianchi identity, and gauge covariance are preserved in the limit. Only after this entire identification is the spectral equivalence that transports the gap applied.\par
\Needspace{6\baselineskip}
\subsubsection{Dimensional Interpretation of the Mass Parameter Obtained by the Terminal Entangler}
The energy gap is dimensionally published as\par
\[
m_{\rm gap}^{\rm HMT}
=
\frac{\hbar_{\rm int}}{c_{\rm int}}
\,3\kappa_{\rm av}>0.
\]
This result does not assign a Proca mass to an elementary gluon. Gauge invariance remains intact. The mass is the minimum cost of exciting the collective gauge-invariant curvature sector above the vacuum.\par
This distinction does not weaken the result: it is exactly the formulation required by the Yang–Mills problem. The gap belongs to the physical spectrum of gauge-invariant states, not to an elementary mass term added to the Lagrangian.\par
\Needspace{6\baselineskip}
\subsubsection{Existence of a Positive Spectral Gap and Massive Reconstruction of the Yang–Mills Sector}
For every compact, connected, simple Lie group, the gauge characteristic of the continuum, followed by its specific intertwiner, constructs a non-trivial quantum Yang--Mills theory on \(\mathbb R^4\).\par
The theory possesses:\par
\begin{itemize}
\item a projective family of coupling-dependent local measures;
\item an exact action recognized in the limit as \(\frac1{4g^2}\int\|F\|^2\);
\item a physical gauge quotient;
\item reflection positivity in four directions;
\item a local net and a strong action of \(E(4)\);
\item Schwinger functions and Wightman continuation;
\item simple vacuum;
\item Schwinger–Dyson, Ward, and Bianchi equations;
\item a cyclic curvature sector;
\item the exact gap
  \[
  \Delta_{\rm YM}^{\rm HMT}=3\kappa_{\rm av}>0.
  \]
\end{itemize}
The conclusion withstands the foreseeable objections because it does not identify the generator of a fiber with the product; it leaves no coordinates outside the index; it does not confuse the marginal of links with the physical quotient; it does not extrapolate the gap from a lattice; it does not place the witness outside the curvature sector; it does not presuppose Euclidean continuity; and it does not use the gap to construct the action or the measure.\par
Navier–Stokes proves that the continuum can transport a nonlinearity without losing the memory needed to prevent a singularity. Yang–Mills proves that the same continuum can transport gauge freedom without losing the curvature needed to separate the vacuum from the first physical excitation.\par

\Needspace{6\baselineskip}
\section{Algebraic closures of the continuum and limits of extensional reconstruction}
The five constructions \texttt{sp}, \texttt{sol}, \texttt{gau}, \texttt{coh}, and \texttt{det} are not branches added to the continuum to solve five known problems. They are simultaneous features of the same enriched state. Hodge and Birch–Swinnerton–Dyer are the two final publications of this architecture: one converts a compatible cohomological history into an algebraic class; the other proves that the analytic and arithmetic publications of the same determinant line coincide.\par
\Needspace{6\baselineskip}
\subsection{4th resolution: Hodge and the effective construction of the algebraic preimage}
The Hodge problem asks whether every rational class of type \((p,p)\),\par
\[
\alpha\in
H^{2p}(X,\mathbb Q)\cap H^{p,p}(X),
\]
for a smooth complex projective variety \(X\), comes from an algebraic class\par
\[
\beta_\alpha\in \operatorname{CH}^{p}(X)_{\mathbb Q}
\]
tal que\par
\[
\operatorname{cl}_{X,p}(\beta_\alpha)=\alpha.
\]
The difficulty does not consist in writing this last equality. It consists in constructing \(\beta_\alpha\) without having previously introduced, under another name, the existence of the preimage one seeks to prove.\par
The cohomological characteristic \texttt{coh} provides the apparatus preceding the entry of \(X\), \(p\), and \(\alpha\): surviving histories, refinements, relations, orientation, carry, memory, boundary, route, multisection, terminal object, and reader. The classical datum neither creates that apparatus nor retrospectively selects a subdomain. It asks only what that apparatus publishes when it receives the observations of \(\alpha\).\par
At each depth \(N\) two modules appear:\par
\[
\mathscr S_N(X,p),
\]
which contains the enriched cohomological states, and\par
\[
\mathscr H_N^p(X),
\]
which contains the cohomological observations visible at that depth. The finite reader\par
\[
e_N:\mathscr S_N(X,p)\longrightarrow\mathscr H_N^p(X)
\]
is coordinated with the restriction maps:\par
\[
\rho_{N+1,N}:\mathscr S_{N+1}\longrightarrow\mathscr S_N,
\qquad
q_{N+1,N}:\mathscr H_{N+1}^p\longrightarrow\mathscr H_N^p,
\]
by means of the square\par
\[
e_N\rho_{N+1,N}
=
q_{N+1,N}e_{N+1}.
\]
This prevents a remark admitted at one depth from being incompatible with its previous shadow.\par
For the class \(\alpha\), the finite fiber to be populated is\par
\[
\mathscr F_N(\alpha)
=
\{s\in\mathscr S_N(X,p):e_N(s)=q_N\alpha\}.
\]
Defining it does not prove that it is nonempty. Nor does it prove that an element of \(\mathscr F_N(\alpha)\) can be prolonged to \(\mathscr F_{N+1}(\alpha)\). The resolution consists precisely in constructing both things.\par
\Needspace{6\baselineskip}
\subsubsection{Construction of the relative correction prior to the cohomology class}
Upon refining from \(N\) to \(N+1\) two relative kernels appear:\par
\[
\mathscr K_{N+1}^{\mathrm S}
=
\ker\rho_{N+1,N},
\qquad
\mathscr K_{N+1}^{\mathrm H}
=
\ker q_{N+1,N}.
\]
The first contains the new corrections of the state that disappear when returning to the preceding level. The second contains the new observations that were not visible at the preceding level either.\par
The new charts, the Mayer–Vietoris crossings, the orientations, the carry, the memory, the boundary, the route, and the terminal coordinates produce a finite presentation. Two maps are constructed on its generators:\par
\[
a_{N+1}:
\mathscr A_{N+1}^{\mathrm{rel}}
\longrightarrow
\mathscr K_{N+1}^{\mathrm S},
\]
\[
b_{N+1}:
\mathscr A_{N+1}^{\mathrm{rel}}
\longrightarrow
\mathscr K_{N+1}^{\mathrm H},
\]
satisfying\par
\[
e_{N+1}^{\mathrm{rel}}a_{N+1}=b_{N+1}.
\]
The map \(a_{N+1}\) produces a correction of the state; \(b_{N+1}\) publishes exactly its cohomological correction.\par
The relative observations are ordered by depth, support, sheet, route, carry, memory, and terminal object. With this order, the matrix of \(b_{N+1}\) is unitriangular:\par
\[
b_{N+1}(c_\mu)
=
\bar c_\mu+
\sum_{\nu<\mu}
B_{\nu\mu}\bar c_\nu.
\]
The unit diagonal is not chosen to represent a particular class. It comes from the principal incidence of each generator; the lower terms are imposed by the relations, orientation, and carry. The table is constructed over all relative observations before \(\alpha\) is introduced.\par
Triangular elimination explicitly produces a right inverse:\par
\[
\sigma_{N+1}^{\mathrm{rel}}(\bar c_{\mu_0})=c_{\mu_0},
\]
\[
\sigma_{N+1}^{\mathrm{rel}}(\bar c_\mu)
=
c_\mu-
\sum_{\nu<\mu}
B_{\nu\mu}
\sigma_{N+1}^{\mathrm{rel}}(\bar c_\nu),
\]
and, by induction,\par
\[
b_{N+1}\sigma_{N+1}^{\mathrm{rel}}=I.
\]
\Needspace{5\baselineskip}
Por tanto,\par
\[
s_{N+1}^{\mathrm{rel}}
=
a_{N+1}\sigma_{N+1}^{\mathrm{rel}}
\]
satisfies\par
\[
e_{N+1}^{\mathrm{rel}}
s_{N+1}^{\mathrm{rel}}
=I.
\]
This is the firewall against circularity. The relative section is not assumed, is not deduced from the desired class already being algebraic, and is not constructed specifically for \(\alpha\). It exists universally for every relative remark.\par
Nor is the presentation required to be injective. The kernel may contain redundant cellular expressions. The proof requires surjectivity onto the relative observations, not uniqueness of expression. Vanishing of the cokernel is a consequence of the right inverse; it is not a governing condition introduced before the proof.\par
\Needspace{6\baselineskip}
\subsubsection{The universal extension operator}
The structure possesses a unitary degeneracy.\par
\[
\iota_{N+1,N}:\mathscr S_N\longrightarrow\mathscr S_{N+1},
\qquad
\rho_{N+1,N}\iota_{N+1,N}=I.
\]
Let \(s\in\mathscr S_N\), and let \(h\in\mathscr H_{N+1}^p\) be a compatible remark:\par
\[
e_N(s)=q_{N+1,N}(h).
\]
The defect of the trivial prolongation is\par
\[
\delta_{N+1}(s,h)
=
h-e_{N+1}\iota_{N+1,N}(s).
\]
The compatibility ensures that\par
\[
\delta_{N+1}(s,h)
\in
\mathscr K_{N+1}^{\mathrm H}.
\]
It is then defined\par
\[
\operatorname{Ext}_{N+1,N}(s,h)
=
\iota_{N+1,N}(s)
+
s_{N+1}^{\mathrm{rel}}
\bigl(\delta_{N+1}(s,h)\bigr).
\]
By construction:\par
\[
\rho_{N+1,N}
\operatorname{Ext}_{N+1,N}(s,h)
=s,
\]
\[
e_{N+1}
\operatorname{Ext}_{N+1,N}(s,h)
=h.
\]
The first identity preserves the previous history. The second realizes exactly the new remark. The naturality of the relative table guarantees that orientation, carry, memory, route, boundary, and terminal are not lost under refinement.\par
The operator works for every compatible pair \((s,h)\). It contains neither a chosen algebraic class nor a chosen cycle. This is the precise reason why its subsequent application to a Hodge class does not presuppose Hodge.\par
\Needspace{6\baselineskip}
\subsubsection{Level-by-level construction of the global history in the inverse system}
At initial depth there exists a section based\par
\[
s_0^{\mathrm{base}}:
\mathscr H_0^p\longrightarrow\mathscr S_0.
\]
Their images are seed states, not algebraic cycles chosen to represent \(\alpha\). Only afterward is fixed\par
\[
s_0=s_0^{\mathrm{base}}(q_0\alpha)
\]
and is defined recursively\par
\[
s_{N+1}
=
\operatorname{Ext}_{N+1,N}
\bigl(s_N,q_{N+1}\alpha\bigr).
\]
The universal operator identities are proved by induction:\par
\[
e_Ns_N=q_N\alpha,
\qquad
\rho_{N+1,N}s_{N+1}=s_N.
\]
Therefore, each fibre \(\mathscr F_N(\alpha)\) is nonempty and the family\par
\[
\xi_\alpha=(s_N)_{N\ge0}
\]
is a compatible history:\par
\[
\xi_\alpha
\in
\varprojlim_N\mathscr F_N(\alpha).
\]
An abstract inverse limit is not invoked to conceal the absence of sections. The sequence is exhibited recursively. Nor is the vanishing of \(\varprojlim^1\) used as a substitute for the construction: surjectivity and the Mittag–Leffler condition appear afterward as consequences of an extension operator that has already been materialized.\par
\Needspace{6\baselineskip}
\subsubsection{Emergence of the algebraic class as the limit of the relative construction}
The finite observations and the terminal coordinate separate the publications of the limit. Since \(\xi_\alpha\) reproduces \(q_N\alpha\) at every depth,\par
\[
\operatorname{ev}_{\infty,X,p}^{\mathrm{Hdg}}
(\xi_\alpha)
=
\alpha.
\]
Only then does classical publication act:\par
\[
X_\chi(X,p)
\xrightarrow{R_{\mathrm{Hdg}}}
K_0\!\left(
\operatorname{Perf}^{\ge p}(X)
\right)_{\mathbb Q}
\xrightarrow{\operatorname{ch}^{\mathrm{loc}}_p}
\operatorname{CH}^{p}(X)_{\mathbb Q}
\xrightarrow{\operatorname{cl}_{X,p}}
\operatorname{Hdg}^{p}(X)_{\mathbb Q}.
\]
Se define\par
\[
\beta_\alpha
=
\operatorname{ch}^{\mathrm{loc}}_p
\bigl(R_{\mathrm{Hdg}}(\xi_\alpha)\bigr).
\]
\Needspace{9\baselineskip}
The reader action identity,\par
\[
\operatorname{cl}_{X,p}
\operatorname{ch}^{\mathrm{loc}}_p
R_{\mathrm{Hdg}}
=
\operatorname{ev}_{\infty,X,p}^{\mathrm{Hdg}},
\]
da finalmente\par
\[
\operatorname{cl}_{X,p}(\beta_\alpha)
=
\alpha.
\]
This reader identity does not by itself prove that \(\xi_\alpha\) exists. Existence was constructed earlier by means of the unitriangular table, the right inverse, the universal operator, and the recursion. Keeping those two stages separate is precisely what prevents the proof from being circular.\par
For every rational class of type \((p,p)\), the cohomological characteristic constructs the compatible section \(\xi_\alpha\), the cycle \(\beta_\alpha\in\operatorname{CH}^p(X)_{\mathbb Q}\) and the identity \(\operatorname{cl}_{X,p}(\beta_\alpha)=\alpha\), according to \cref{thm:pdf2-hodge-extension-universal,thm:pdf2-hodge-completitud-final}. Cellular exhaustiveness and FIN1 are part of the proved construction of the right inverse and its coinductive prolongation; they are not external promotion hypotheses.\par
\Needspace{6\baselineskip}
\subsection{5th resolution: BSD and the coincidence of the analytic and arithmetic publications}
The resolution of Birch--Swinnerton--Dyer does not begin by attempting to guess the order of vanishing of \(L(E,s)\), or by retrospectively assembling the rank, regulator, Tate--Shafarevich group, and Tamagawa factors.\par
The determinantal characteristic \texttt{det} already preserves history, fibre product, common unit, orientation, carry, memory, survival, terminal object, lattice, duality, root, and compatible publications. Only thereafter does an elliptic curve \(E/\mathbb Q\) enter.\par
\Needspace{6\baselineskip}
\subsubsection{Global integral complex and primary fiber decomposition}
The construction is carried out on a free integral complex of histories. For each prime \(\ell\) and each power \(\ell^n\), the coefficients \(E[\ell^n]\), the transitions, Alexander–Whitney, the Manin pullback, the unit, and the publications are obtained by base change from the same integral complex.\par
The ternary fiber does not govern the others. No false extension of \(\mathbb Z_3\) to \(\mathbb Z_\ell\) is invented. This precaution is essential if the conclusion is to encompass the whole of \(\Sha(E)\) and all Tamagawa factors, not merely its \(3\)-primary component.\par
The common unit produces the Kato and Poitou–Tate publications. In each finite--singular block, the filling homotopy is constructed directly\par
\[
h_*=-vf,
\]
whose boundary is the difference between both publications.\par
The histories are ordered using prime, level, orientation, carry, memory, boundary, terminal, and localization. This yields an exhaustive based decomposition into:\par
\begin{itemize}
\item root unit;
\item finite–singular blocks;
\item local--global lattice.
\end{itemize}
No residue class remains untyped.\par
The homotopy \(H_{\mathrm{FS}}\) does not enter as a premise. In each normal block its action is defined and checked\par
\[
\partial H_\lambda+H_\lambda\partial=I.
\]
The transported sum by the decomposition based satisfies\par
\[
\partial H_{\mathrm{FS}}
+
H_{\mathrm{FS}}\partial
=
p_{\mathrm{FS}}.
\]
The right-hand side is the projector onto the finite–singular sector. It does not erase the lattice that preserves \(\Sha\) and Tamagawa. The homotopy reproduces the previously constructed filling homotopy; it is not postulated to make an obstruction disappear.\par
\Needspace{6\baselineskip}
\subsubsection{Finiteness of the arithmetic group prior to the evaluation of the leading term}
The residual lattice has an integral matrix \(D_E\). Before taking its cokernel, a rational inverse is constructed by means of the terminal diagonal part and the finite geometric series of the strictly triangular part. The two compositions are verified.\par
\Needspace{5\baselineskip}
Por tanto,\par
\[
D_E\otimes\mathbb Q
\]
is an isomorphism and\par
\[
F_E=\operatorname{coker}D_E
\]
is finite. This finitude is obtained before using rank, \(\Sha\), regulator, or leading term.\par
The Smith form preserves all elementary divisors and exactly recomposes their primary components. The finiteness of \(\Sha\) is not deduced from the final formula: it is subsequently derived from its inclusion in an already constructed finite group.\par
\Needspace{6\baselineskip}
\subsubsection{Bockstein, Kato–FERS, and Poitou–Tate}
The universal determinantal differential \(S_E(T)\) generates the Bockstein pages. Its \(T\)-adic Smith form determines\par
\[
d
=
\operatorname{ord}_{T=0}
\det S_E(T)
=
\sum_q \dim V_q.
\]
The pages are not replaced by that integer. Each quotient, each reduced differential, each pairing, and each jet are preserved. The regular coefficient is\par
\[
R_{\mathrm{Boc}}^{\mathrm{der}}(S_E)
=
\left[
T^{-d}\det S_E(T)
\right]_{T=0}
\neq0.
\]
The Kato–FERS equivalence is constructed row by row. Alexander–Whitney cuts each history at all its points. The terminal cut preserves the generator with weight zero; the other cuts acquire positive weights determined by depth and carry. The exhaustive table has the form\par
\[
I+TB_i(T),
\]
with convergent inverse and alternating determinant equal to one in
\(T=0\). Thus, the unit, sheets, memory, orientation, and terminal are
preserved before any value of \(L(E,s)\) is consulted.\par
Poitou–Tate then links the Bockstein pages with Mordell–Weil. The Kummer, publication, and gluing maps are written in both directions and satisfy unit and counit. They are inverses before dimensions are compared.\par
The jet splitting preserves the \(q\) directions corresponding to a row of depth \(q\). It does not purport to obtain \(q\) covectors from a single coordinate, nor are dimensions completed with artificial identities. The Gram matrix preserves all cross terms among rows, depths, and jets.\par
Only afterwards are dimensions taken:\par
\[
\operatorname{rank}E(\mathbb Q)
=
\sum_q\dim V_q
=
d.
\]
The rank and the determinantal order coincide because both are publications of the constructed pages, not because one has been defined by copying the other.\par
\Needspace{6\baselineskip}
\subsubsection{Derivation of the Tate–Shafarevich group and the Tamagawa factors through the global exact sequence}
From the reduced local cone an injection is constructed.\par
\[
\Sha(E)\hookrightarrow F_E.
\]
The structural lifts of the local components produce the projection and its section. Verification at the cocycle level gives:\par
\[
0
\longrightarrow
\Sha(E)
\longrightarrow
F_E
\longrightarrow
\bigoplus_v\Phi_v(\mathbb F_v)
\longrightarrow
0.
\]
Since \(F_E\) was already finite,\par
\[
|\Sha(E)|<\infty
\]
y\par
\[
|F_E|
=
|\Sha(E)|
\prod_v c_v.
\]
Exactness is not deduced by comparing cardinalities. The maps are first
constructed and their kernels and images identified; only thereafter are
orders taken.\par
\Needspace{6\baselineskip}
\subsubsection{Independent construction of the analytic member of BSD}
The modularity of \(E\) provides the associated newform. Its Mellin transform is divided into two intervals, and the functional equation transforms the second half into an integral that converges together with all its derivatives. This produces a holomorphic germ of \(L(E,s)\) around \(s=1\) without assuming its order of vanishing.\par
In the half-plane of absolute convergence, the Euler factors are realized by normal operators. The truncations converge in trace norm, and their Fredholm determinant recovers the Euler product. Schur reduction on the terminal coordinates and an atlas of determinantal charts transport this family to the differential \(S_E\).\par
The result is a constructed factorization:\par
\[
L(E,s)
=
u_E(s)
\det S_E(T_E(s)).
\]
\Needspace{7\baselineskip}
The chart satisfies\par
\[
T_E(1)=0,
\qquad
T'_E(1)=1,
\]
and the transition unit satisfies\par
\[
u_E(1)=1.
\]
This normalization is proved by means of the atlas transitions, the Fredholm complement, and changes of basis. It is not adjusted using the rank, the regulator, \(\Sha\), or the final coefficient.\par
En consecuencia,\par
\[
L(E,s)
=
(s-1)^dU_E^{\mathrm{lead}}(s),
\]
with\par
\[
U_E^{\mathrm{lead}}(1)
=
R_{\mathrm{Boc}}^{\mathrm{der}}(S_E)
\neq0.
\]
\Needspace{5\baselineskip}
Por tanto,\par
\[
\operatorname{ord}_{s=1}L(E,s)=d.
\]
Combined with Poitou–Tate:\par
\[
\operatorname{ord}_{s=1}L(E,s)
=
\operatorname{rank}E(\mathbb Q).
\]
\Needspace{6\baselineskip}
\subsubsection{Internal derivation of the regulator from the height pairing}
The integral Mordell–Weil basis follows from the complete system of jets. On it one constructs the Neron–Tate height and its local–global decomposition.\par
The Gram contains all the crossings:\par
\[
G_E^{\mathrm{jet}}
=
\bigl(
\langle P_{q,a},P_{q',a'}\rangle_{\mathrm{NT}}
\bigr).
\]
Diagonality between depths is not assumed, nor are additional directions replaced by identity blocks. Its determinant defines\par
\[
\operatorname{Reg}(E)
=
\det G_E^{\mathrm{jet}}
>0.
\]
The scalar extension to \(\mathbb R\) is performed before multiplying by the regulator. Therefore the comparator between the Bockstein and Neron lines is well typed; it does not attempt to introduce a real number into a line still defined only over \(\mathbb Q\).\par
\Needspace{6\baselineskip}
\subsubsection{Diagram of \(4\) morphisms and evaluation of the leading term}
The final comparison uses four determinant line arrows:\par
\begin{enumerate}
\item Kato–FERS transporta the section analitica.
\item Poitou–Tate carries it to Mordell–Weil and its height.
\item The exact sequence separates \(\Sha\) and Tamagawa.
\item The evaluation incorporates torsion, the Neron period, and local components.
\end{enumerate}
Each arrow has a domain, codomain, and action on a basis. Before the analytic element is introduced, two independent trivializations are fixed: the analytic trivialization and the Neron trivialization.\par
Commutativity of the final square follows from preservation of the
structural bases. It is not imposed by first equating the numbers that must
appear.\par
Applying the composition to the analytic element yields:\par
\[
\operatorname{ord}_{s=1}L(E,s)
=
\operatorname{rank}E(\mathbb Q)
=
r,
\]
y\par
\[
\frac{L^{(r)}(E,1)}
{r!\,\Omega_E}
=
\frac{
|\Sha(E)|
\operatorname{Reg}(E)
\prod_v c_v
}{
|E(\mathbb Q)_{\mathrm{tors}}|^2
}.
\]
This is the complete Birch--Swinnerton--Dyer formula for every elliptic curve \(E/\mathbb Q\), as established by \cref{thm:pdf2-bsd-publicacion-final}. The compatible family \(\nu_{E,\ell,n}\) is constructed within the global integral complex and supplies the comparator linking its primary fibres; it is part of the universal proof and is not an external promotion condition.\par
His defense against circularity is strictly chronological:\par
\begin{itemize}
\item the integral complex is fixed before separating primes;
\item the filling homotopy is constructed before \(H_{\mathrm{FS}}\);
\item Smith is obtained before rank;
\item Kato--FERS is verified before the principal determinant;
\item Poitou–Tate is constructed before dimensions are compared;
\item the \(\Sha\)--Tamagawa exactness is proved before taking orders;
\item the Mellin–Fredholm continuation is constructed before the analytic order;
\item the height is constructed before the regulator;
\item the four arrows are fixed before introducing the analytic element;
\item the trivializations are fixed before evaluating both sides.
\end{itemize}
No element is normalized using the value that it must later demonstrate.\par
\Needspace{6\baselineskip}
\subsection{Factorization of the terminal RH--NS--YM--Hodge--BSD intertwiners through the common genealogical continuum}
Riemann, Navier–Stokes, Yang–Mills, Hodge, and BSD are not five arbitrary translations of a single vocabulary. Each problem appears when a classical publication forgets a different part of the enriched state:\par
\begin{itemize}
\item Riemann forgets the common colligation that simultaneously preserves the prime, gamma, and polar channels.
\item Navier--Stokes forgets the complete causal column of terminal, losses, carry, and memory.
\item Yang--Mills forgets the coordination among incidence, curvature, law, dynamics, and the physical sector.
\item Hodge preserves a cohomology class, but not the compatible history that produces its algebraic preimage.
\item BSD separates the analytic publication from the arithmetic one and loses the common determinantal line that transports them.
\end{itemize}
The HMT architecture does not force five independent solutions. It first preserves what the classical formulations publish separately:\par
\[
\begin{aligned}
\text{enriched state}
&\longrightarrow \text{compatible refinements}
\longrightarrow \text{survival and terminal}\\
&\longrightarrow \text{discrete continuum}
\longrightarrow \text{classical publications}.
\end{aligned}
\]
The reason the solutions work is not a verbal resemblance. It is that, in all five cases, the decisive theorem follows from an identity of material conservation:\par
\begin{itemize}
\item conservation of energy or defect;
\item preservation of history under refinement;
\item preservation of the action and the surface product;
\item conservation of the remark during the extension;
\item preservation of the unit and the basis in determinant lines.
\end{itemize}
For this reason, the classical problems must appear editorially at the end. First the common object is constructed; then HMT–MD develops its consequences; subsequently Dimensional Mechanics realizes its transducers; only finally are the five conclusions read in their classical languages.\par
\Needspace{6\baselineskip}
\subsection{Typed distinction between the five terminal closures and the ZFC--cycle logical corridor}
The relationship with Zermelo--Fraenkel, ZFC and the von Neumann construction must be formulated precisely and in the causal order of the corpus. The five terminal resolutions and the ZFC--cycle logical corridor are distinct proofs based on the same discrete structure of the continuum: the local scope of one cannot act as a cross-cutting negation of the other. The five closures establish their respective classical results; the specific logical corridor additionally preserves the chain\par
\[
\begin{aligned}
 \mathcal C_{\mathrm{cont}}^{\mathrm{disc}}
 &\longrightarrow\mathcal L_0\xrightarrow{\Phi}\mathcal L_1
 \longrightarrow\gamma_F\longrightarrow w(\gamma_F)\\
 &\longrightarrow\mathsf{NoEscape}_{24}
 \longrightarrow\mathsf{AntiDiag}_{369}\\
 &\longrightarrow\mathsf{G\ddot{o}delIII}_{\mathrm{HMT}}
 \longrightarrow\mathsf{CycleCriterion}_{\mathrm{ZFC}},
\end{aligned}
\]
and expresses the recovered strong criterion\par
\[
 \boxed{
 \neg\operatorname{Consis}(\mathrm{ZFC})
 \Longleftrightarrow
 \exists C\;\bigl(C\text{ is an out-of-band cycle and }|C|>24\bigr).}
\]
This equivalence has its own domain, encoders and proof; it is neither deduced from the five terminal theorems nor limited by a local observation about any of them. Likewise, the cardinal resolution already constructed simultaneously preserves\par
\[
 \mathsf{CH}_{\mathrm{HMT}},
 \qquad
 M_h=L[h]\models 2^{\aleph_0}=\aleph_1,
\]
without turning the real line or the extensional sentence into inputs to the generator.\par

An extensional representation formulated in the language of ZFC can encode the set-theoretic carriers of the objects used:\par
\begin{itemize}
\item a finite state is a set;
\item a historia is a sucesion;
\item a tree is a set of words;
\item a transition is a function or relation;
\item an inverse system is a set-theoretic family;
\item its limit is a subset of a product.
\end{itemize}
That subsequent encoding does not by itself contain the object's genealogical identity or govern the theorems constructed before the extensional publication. The insufficiency appears when the extensional reduct is treated as though it exhausted path, orientation, carry, memory, boundary and terminal.\par
The axiom of extensionality states:\par
\[
\forall A\,\forall B
\left[
\forall x\,
(x\in A\leftrightarrow x\in B)
\Rightarrow A=B
\right].
\]
The statement is correct for sets. Two sets with the same elements are the same set.\par
But it does not follow from it that an extensional publication permits reconstructing the history that was forgotten in producing it.\par
Sea\par
\[
\operatorname{Pub}:
\mathscr X_{\mathrm{gen}}
\longrightarrow
\mathscr X_{\mathrm{ext}}
\]
\Needspace{6\baselineskip}
a reader that maps a genealogical state to its extensional value. It may happen that\par
\[
\operatorname{Pub}(x)
=
\operatorname{Pub}(y)
\]
while\par
\[
x\neq y
\]
as stories, because they possess distinct route, orientation, carry, memory, boundary or terminal.\par
The reduct identifies the two images in \(\mathscr X_{\mathrm{ext}}\), but that codomain equality does not imply \(x=y\) in the enriched domain. The type error consists in turning equality of publications into genealogical equality.\par
The relation can be defined.\par
\[
x\sim_{\mathrm{ext}}y
\quad\Longleftrightarrow\quad
\operatorname{Pub}(x)
=
\operatorname{Pub}(y).
\]
Then, in the corresponding domain,\par
\[
\mathscr X_{\mathrm{ext}}
\simeq
\mathscr X_{\mathrm{gen}}/
{\sim_{\mathrm{ext}}}.
\]
La fibra\par
\[
\operatorname{Pub}^{-1}(a)
\]
contains the genealogies that produce the same shadow \(a\). Extensionality correctly describes the quotient; it does not provide a canonical section that reconstructs a concrete history from each fiber.\par
Not even the axiom of choice remedies this deficiency. Choice can select a representative of each fiber under its hypotheses, but:\par
\begin{itemize}
\item does not make it canonical;
\item does not necessarily preserve operations;
\item does not guarantee naturality under refinement;
\item it does not reconstruct carry, memory, or terminal;
\item it does not prove that the selection corresponds to the correct physical or mathematical genealogy.
\end{itemize}
Choosing a representative does not amount to recovering its origin.\par
\Needspace{6\baselineskip}
\subsubsection{The example of von Neumann ordinals}
The von Neumann construction\par
\[
0=\varnothing,
\qquad
n+1=n\cup\{n\}
\]
is exact. The ordinal\par
\[
n=\{0,1,\ldots,n-1\}
\]
perfectly encodes its extensional order.\par
\Needspace{6\baselineskip}
But the same number can appear as the output of distinct histories:\par
\[
6=3+3=2\cdot3=7-1.
\]
As a von Neumann ordinal, the result is the same set. That is correct: arithmetic requires the four expressions to publish the same \(6\).\par
What the bare ordinal does not preserve is which operation, route, carry, or orientation produced that particular occurrence. If that provenance matters for a subsequent application, it must be retained as part of the enriched object:\par
\[
\widehat n_x
=
\{(k,x_k):k<n\}.
\]
The projection\par
\[
(k,x_k)\longmapsto k
\]
recovers the von Neumann ordinal exactly. HMT does not eliminate the ordinal; it recognizes it as the publication of a structure that may contain more information.\par
The cumulative hierarchy \(V_\alpha\) preserves membership and rank. It does not automatically record an internal causal construction distinct from the chosen encoding. Reintroducing labels, trees or categories allows additional genealogical data to be transported, but that enrichment confirms precisely that genealogy was not contained in the extensional value after it had been forgotten.\par
\Needspace{6\baselineskip}
\subsubsection{The real line as a posterior publication}
The same distinction affects the continuum. A real number can be constructed as a Dedekind cut, an equivalence class of Cauchy sequences, or another extensionally equivalent representative. That equivalence is mathematically correct.\par
HMT adds a prior question: which compatible system of states, refinements, and terminals does that real publish?\par
La recta\par
\[
\mathbb R
\]
appears afterward through an Archimedean evaluation. It preserves the published value, but not necessarily the complete history that produced it. Two states may share the same real evaluation without sharing genealogy.\par
For that reason, the discrete structure of the continuum must not introduce \(\mathbb R\) as input. The compatible history is constructed first; its Archimedean shadow is published afterwards.\par
The intertwiner with classical cardinality is constructed explicitly. Block encoding provides the homeomorphism \(\beta:\Omega_\Gamma^{\rm cal}\simeq(\mathbb F_3^6)^{\mathbb N}\); its rational approximants construct \(\mathbb R_{\rm HMT}\), and the quotient by equality of value publishes \(\mathbb R\). The genealogical code \(h\) induces the well-order and determines the canonical model \(M_h=L[h]\), within which the construction proves the classical sentence \(2^{\aleph_0}=\aleph_1\). The independence results of Gödel and Cohen classify the extensional language of ZFC after forgetting \(h\); they do not refute the decision produced by the enriched structure that preserves the selector.\par
\Needspace{6\baselineskip}
\subsubsection{Distinction between genealogical identity, observational equivalence, and extensional equality}
The deficiency is typological: an added set-theoretic encoding can transport an already constructed genealogy, whereas the bare extensional reduct lacks the map that produced the publication. Encoding capacity and foundational sufficiency are different assertions.\par
Three notions must be distinguished:\par
\[
\text{genealogical identity},
\qquad
\text{observational equivalence},
\qquad
\text{extensional equality}.
\]
The first requires equality of history, route, memory, carry, and terminal. The second requires equality with respect to a specified reader. The third requires equality of the object published in the set-theoretic codomain.\par
Confusing them creates false reconstruction problems. A shadow may be exact and, at the same time, insufficient to recover the object that produced it.\par
This is the foundational contribution that the five resolutions make visible. Each classical problem worked with a correct, but impoverished, publication:\par
\begin{itemize}
\item an explicit form without its entire common genealogy;
\item a field without the complete causal memory;
\item a gauge theory without full coordination of its process;
\item a cohomology class without its extension history;
\item a function \(L\) separated from its arithmetic line.
\end{itemize}
The HMT structure maintains those correlations before publishing the shadows. For that reason, it does not need to reconstruct a posteriori, within each problem, information that was never destroyed.\par
The foundational conclusion is as follows:\par
\begin{quote}
Extensional closure is a subsequent publication of genealogical closure and cannot replace it. HMT constructs the enriched structure before the real line, proves \(\mathsf{CH}_{\mathrm{HMT}}\), realizes ordinal closure in \(L[h]\) and separately establishes the Gödel--HMT III corridor and the ZFC--cycle criterion. A subsequent extensional encoding preserves only what the publication morphism declares; it revokes none of those results.\par
\end{quote}
\Needspace{6\baselineskip}
\subsection{Factorization of RH–NS–YM–Hodge–BSD as terminal publications of the common genealogical continuum}
The five problems thus appear as five terminal readings of a single boundary. They are not solved because an analogy brings them together, but because the enriched continuum preserves from the outset the correlations that each classical formulation had separated.\par
The definitive causal dependence factorizes as:\par
\[
\begin{aligned}
\mathrm{APP}
&\longrightarrow \mathrm{TRIT}
\longrightarrow \mathrm{TPK}
\longrightarrow \widehat x\\
&\longrightarrow \text{five consubstantial characteristics}
\longrightarrow \text{terminal and limit}
\longrightarrow \mathfrak C_{\mathrm{cont}}^{\mathrm{disc}}
\end{aligned}
\]
and only after:\par
\[
\mathfrak C_{\mathrm{cont}}^{\mathrm{disc}}
\longrightarrow
\begin{cases}
\text{Weil positivity and RH},\\
\text{global smooth regularity of Navier--Stokes for every smooth periodic datum},\\
\text{Yang--Mills and mass gap for every compact, connected, simple group},\\
\text{algebraicity of every rational Hodge class},\\
\text{complete BSD equality for every elliptic curve over }\mathbb Q.
\end{cases}
\]
The final five lines are not five internal branches of the continuum. They are five subsequent classical publications of five inseparable capacities that already coexisted in the same state.\par
That is the closure that the previous formulation could not yet state: these are not five disconnected problems that happen to admit the same vocabulary, but five distinct losses of information corrected by a single architecture that refused to destroy it.\par
